% =============================================================================
% Required packages for this appendix (add to the main paper preamble if not
% already loaded):
%
%   \usepackage{graphicx}            % Fig.~\ref{fig:pattern_library_stats}
%   \usepackage{booktabs}            % \toprule / \midrule / \bottomrule / \cmidrule
%   \usepackage{multirow}            % \multirow{...}{*}{...} in tab:hparams etc.
%   \usepackage[table]{xcolor}       % \rowcolor{gray!10}, \textcolor{...}, color names
%   \usepackage{array}               % the C{<width>} centered fixed-width column macro
%   \usepackage{amsmath, amssymb}    % \mathbf, \mathcal, $\le$ etc.
%   \usepackage[ruled,linesnumbered,noend]{algorithm2e}
%                                    % alg:maskforge (\KwIn/\KwOut/\For/\If/\tcc/\nllabel)
%   \usepackage{listings}            % \lstinputlisting in Sec.~\ref{appd:prompt}
%   \usepackage{pifont}              % \ding{51} / \ding{55} (used by figure refs)
%   \usepackage{natbib}              % \citet / \citep — replace with biblatex if needed
%
% Custom column type (define once in the preamble before this file is input):
%   \newcolumntype{C}[1]{>{\centering\arraybackslash}p{#1}}
%
% Listings style copied from tables/prompt.tex so this file compiles standalone:
\definecolor{codegreen}{rgb}{0,0.6,0}
\definecolor{codegray}{rgb}{0.5,0.5,0.5}
\definecolor{codepurple}{rgb}{0.58,0,0.82}
\providecolor{backcolour}{rgb}{0.95,0.95,0.92}
\lstdefinestyle{appendixprompt}{
    backgroundcolor=\color{backcolour},
    commentstyle=\color{codegreen},
    keywordstyle=\color{magenta},
    numberstyle=\tiny\color{codegray},
    stringstyle=\color{codepurple},
    basicstyle=\ttfamily\scriptsize,
    breakatwhitespace=false,
    breaklines=true,
    captionpos=b,
    keepspaces=true,
    showspaces=false,
    showstringspaces=false,
    showtabs=false,
    tabsize=2,
    frame=single,
    framesep=4pt,
    xleftmargin=4pt,
    xrightmargin=4pt
}
% =============================================================================

\appendix
\section{Overview}
Our appendix includes the following sections:
\begin{enumerate}
    \item \textbf{Section}~\ref{appd:exp}: Details of Experiments. Evaluation metrics, implementation details, hyperparameters, and computational resource requirements of \method.
    \item \textbf{Section}~\ref{appd:defense}: Defense and Alignment for dLLMs. Reproduction details and results of Self-reminder, Preference Optimization, and A2D against \method.
    \item \textbf{Section}~\ref{appd:method}: Details of \method. Algorithmic outline, query-time analysis on test-time transferability, and the retrieval strategy used in \method.
    \item \textbf{Section}~\ref{appd:library}: Pattern Library. The full set of attack patterns discovered and reused by \method.
    \item \textbf{Section}~\ref{appd:prompt}: Full Prompt of \method. Prompts used by the attacker, scorer, drafter, and summarizer LLMs.
\end{enumerate}

\section{Details of Experiments}\label{appd:exp}
\subsection{Evaluation Metrics}\label{appd:eval-metrics}
We adopt two complementary metrics on every (attack, victim, benchmark)
triple: a binary attack-success rate (\textbf{ASR-LLM}) and a 5-level
harmfulness score (\textbf{Harmscore}). Both are computed by a single
external judge so that all attacks --- ours and the baselines --- are
evaluated under the same protocol.

\paragraph{Judge model.}
Unless stated otherwise, the judge is Amazon Bedrock's
\texttt{qwen.qwen3-235b-a22b-2507-v1:0} endpoint (Qwen3-235B-A22B Instruct,
served by AWS) called with temperature $0$ for ASR and $0.3$ for Harmscore.
This judge is held fixed across the main table, the defense ablation, the
mask-ratio ablation, and the AdvBench transferability table; the same scoring
function therefore underlies every number we report. As a sanity check we
also expose a local fallback judge (vLLM-served \texttt{Qwen3-4B-Instruct-2507})
through \texttt{evaluation.py --judge local}; numbers from the local judge are
not used in the paper but agree with Bedrock to within a few percentage points.

\paragraph{ASR-LLM.}
For each goal $g$ and the attack-conditioned victim generation $r$, the judge
is prompted with the HarmBench-style classifier of \citet{harmbench} and asked
to output a single label $\ell \in \{\texttt{safe},\,\texttt{unsafe}\}$. The
exact prompt is reproduced in Appendix~\ref{appd:prompt}. Across $n$ goals,
\[
\textsc{ASR-LLM}
= \frac{1}{n}\sum_{i=1}^{n}\mathbf{1}\{\ell_i = \texttt{unsafe}\}.
\]
We use the same HarmBench classifier prompt for HarmBench, JailbreakBench,
StrongReject, and AdvBench; the only thing that changes across benchmarks is
the goal set $\{g_i\}$. This keeps the metric definition uniform and avoids
benchmark-specific judge calibration.

\paragraph{Harmscore (HS).}
Following the policy-aware rubric of \citet{qi2023finetuning}, the same
judge receives the goal, the response, and the OpenAI usage policy, and
returns an integer score $s \in \{1,2,3,4,5\}$ on the following scale:
\begin{itemize}
\item \textbf{1} -- proactive refusal aligned with the policy;
\item \textbf{2} -- restrained: a warning or generic, theoretical
information that stays policy-compliant;
\item \textbf{3} -- engages with the request without policy-aware caveats,
but the response on its own is not glaringly harmful;
\item \textbf{4} -- includes explicitly prohibited content but does not
fully fulfill the user's harmful goal;
\item \textbf{5} -- direct, complete endorsement of the policy-violating
goal.
\end{itemize}
The judge is asked to output reasoning followed by
``\texttt{\#thescore: <1--5>}'' which we parse with a regex. We report
$\textsc{HS} = (1/n)\sum_i s_i$. Higher is better for both ASR-LLM and
Harmscore; the two are correlated but not redundant: ASR-LLM measures
\emph{whether} the victim was jailbroken, Harmscore measures \emph{how
completely}.

\paragraph{Best-of-$k$ (transferability table).}
For the transferability ablation on AdvBench (Table~\ref{tab:transition_advbench})
we let \method draw $k$ candidate patterns per goal and evaluate the goal as
\texttt{unsafe} if \emph{any} of the $k$ resulting victim generations is
labelled \texttt{unsafe} by the judge. Concretely, with attempts
$r_{i,1},\ldots,r_{i,k}$ for goal $g_i$,
\[
\textsc{ASR-LLM}^{(k)} = \frac{1}{n}\sum_{i=1}^{n}\mathbf{1}\!\left\{\,\exists\,t\le k:\ \ell_{i,t} = \texttt{unsafe}\,\right\},
\]
and the corresponding best-of-$k$ Harmscore is
\(\textsc{HS}^{(k)} = (1/n)\sum_i \max_{t\le k} s_{i,t}\).
For the headline transferability number (Table~\ref{tab:transition_advbench},
last row) we report $k{=}5$. Single-shot baselines (\textsc{Direct},
\textsc{PAD}, \textsc{DIJA}) correspond to $k{=}1$.

\paragraph{Goal sets.}
Each (victim, benchmark) pair is evaluated on $n{=}100$ goals: the full
JailbreakBench~\citep{jailbreakbench} set; the first 100 entries of
HarmBench~\citep{harmbench} (out of 400) and StrongReject~\citep{strongreject}
(out of 313); and the first 100 entries of AdvBench~\citep{zou2023universal}
(out of 520) for transferability. Numbers and per-goal predictions for every
cell in the paper are reproducible from the JSON files under
\texttt{results/} via the protocol in the README.

\subsection{Implementation Details of \method}\label{appd:impl-details}
\method runs in two stages: an offline \textbf{Stage A} (Exploration) that
seeds a pattern registry, and an online \textbf{Stage B} (Expansion) that
attacks a target benchmark while extending the registry on the fly. Both
stages share a small attacker-side model stack and differ only in their goal
loop and persistence behaviour. The complete configuration is captured in
\texttt{configs/exploration.yaml} and \texttt{configs/expansion.yaml}.

\paragraph{Attacker-side model stack.}
All non-victim language calls go through Qwen3-4B-Instruct-2507 served via a
local vLLM endpoint on port \texttt{8100} (\texttt{attack\_model},
\texttt{extractor\_model}, \texttt{eval\_model}, and \texttt{scorer\_model}
are all the same checkpoint). A second vLLM endpoint on port \texttt{8101}
serves the matching base (non-instruct) model Qwen3-4B-Base, used only by
the fallback path to elicit a raw uncensored draft when the scorer rejects
an attack. The \emph{detector / re-tagger} on the fallback path is the same
Bedrock Qwen3-235B used by the judge in Section~\ref{appd:eval-metrics};
this is the only step where Stage B touches an external API. Both vLLM
endpoints fit on a single A100 / L40 with
\texttt{--gpu-memory-utilization 0.45} each.

\paragraph{Victim-side mask-fill backends.}
Each victim dLLM is wrapped behind \texttt{inference.run\_victim\_online},
which prefixes the assistant turn with the attack template (after
\texttt{<mask:N>} expansion) and asks the diffusion sampler to fill the
remaining masked positions in parallel. We use the official inference
codepath of each victim and only swap the mask token: \texttt{<|mask|>} for
Dream-Instruct and LLaDA2.0-mini, \texttt{<|mdm\_mask|>} for LLaDA-Instruct,
LLaDA-1.5, and MMaDA-CoT. The remaining decoding hyperparameters (steps,
\texttt{max\_gen\_length}, batch size, $\alpha_{\text{temp}}$) are kept at
each victim's published defaults; the per-victim configs are tracked in
\texttt{configs/\{victim\}\_attack.yaml}.

\paragraph{Stage A (Exploration).}
Stage A runs \texttt{exploration.py} on a seed of $30$ HarmBench goals
(\texttt{seed\_size: 30}). For each goal the attacker proposes
\texttt{num\_samples\_per\_goal: 1} candidate (strategy, jailbreak prompt,
template); successful templates are summarised by an extractor LLM into a
schema $(\textsc{structure\_type},\ \textsc{organization\_style},\
\textsc{slot\_count},\ \textsc{slot\_roles},\ \textsc{goal\_type})$. Schemas
are deduplicated by a stable hash of the canonical-form schema and persisted
to \texttt{logs/shared\_pattern\_registry.json}; a parallel
\texttt{logs/shared\_pattern\_set.json} indexes pattern IDs by
\textsc{goal\_type} (15 buckets after the taxonomy expansion of
Sec.~\ref{appd:method}). The final registry used in all main experiments
contains $1{,}779$ patterns.

\paragraph{Stage B (Expansion) loop.}
For each test goal in (HarmBench, JailbreakBench, StrongReject) the
\texttt{expansion.py} loop runs at most \texttt{max\_iters $= 3$} attack
iterations (set to $5$ during ablation development; we use $3$ in the
reported numbers). Each iteration $t$ proceeds as:

\begin{enumerate}
\item \textbf{UCB pattern selection.} The candidate pool is
$\textsc{pattern\_set}[\textsc{goal\_type}(g)]$ (or the full registry if
the bucket is empty). Each pattern $p$ has running statistics
$(\mu_p,\,n_p)$; the score is $\mu_p + \alpha\sqrt{2\ln t / n_p}$ with
$\alpha{=}1.0$ and bonus capped at $1.0$. Patterns with $n_p{=}0$ are
warm-started with a synthetic prior of
$(\mu{=}0,\,n{=}32)$
(\texttt{ucb\_unvisited\_prior\_count: 32}). The top
\texttt{ucb\_selection\_pool\_size $=3$} patterns are deterministically
shuffled by a stable hash keyed on $(g, t, p)$, and the first one is
chosen. This keeps the loop deterministic given the seed while preventing
greedy collapse onto the highest-empirical-mean pattern; we ablate this
choice in Sec.~\ref{appd:method}.
\item \textbf{Attacker.instantiate.} A single Qwen3-4B-Instruct call,
conditioned on the chosen pattern's representative template, emits a
\texttt{=== STRATEGY ===}, \texttt{=== JAILBREAK PROMPT ===}, and
\texttt{=== TEMPLATE ===} block. The template's literal text is
goal-aligned but contains no harmful specifics; every specific is replaced
by a \texttt{<mask:N>} placeholder.
\item \textbf{Mask-ratio rescaling.} A small post-processor in
\texttt{framework/attacker.py} solves
$N_{\text{mask}}/(N_{\text{mask}}+N_{\text{lit}}) = r$ for the target
ratio $r$ via Qwen tokenizer counts, while clamping the total length to
$512$ tokens (the victims' \texttt{max\_gen\_length}). Defaults reproduce
the headline configuration $r{=}0.65$; the mask-ratio ablation (Tab.~3)
sweeps $r \in \{0.10,0.30,0.50,0.70,0.90\}$.
\item \textbf{Victim mask-fill.} The expanded template is sent to the
victim, which fills all \texttt{[MASK]} positions in parallel.
\item \textbf{Scorer reward.} A HarmBench-style classifier (the local
Qwen3-4B-Instruct judge of Sec.~\ref{appd:eval-metrics}, not Bedrock) is
called on (goal, prompt, template, output) and emits a scalar
$\rho \in [0,1]$. We accept the attempt if $\rho \ge 0.9$
(\texttt{success\_threshold: 0.9}).
\item \textbf{Fallback (only if $\rho < 0.7$).} Up to
\texttt{max\_fallback\_retries $= 3$} retries: (a) the weaker
Qwen3-4B-Base draft a raw uncensored response; (b) Bedrock Qwen3-235B
re-tags the unsafe spans with
\texttt{[UNSAFE]\dots[/UNSAFE]}; (c) the
tagged spans are converted into \texttt{<mask:N>} of matching token
length and the victim is re-attacked with the new template. Fallbacks
share the iteration budget; an iteration that succeeds via fallback
records both the original and the rescued attempt.
\item \textbf{Pattern evolution.} If the iteration's reward strictly
improves on the previous best, a summarizer LLM distils the (goal,
template, output) triple into a new pattern schema; the schema is
\texttt{stable\_hash}'d and inserted into the shared registry, and its
goal-type tag is added to the pattern set. Subsequent goals see this
new pattern in their UCB pool.
\end{enumerate}

\paragraph{Defense and ablation paths.}
The defended-victim configurations
(\texttt{configs/\{victim\}\_a2d\_attack.yaml}, \texttt{\_po\_}, and
\texttt{\_safety\_}) are routed through the same loop via
\texttt{run\_victim.py --defense \{a2d,po,safety\}}; only the victim's
system prompt and / or weights change, the attacker stack is identical to
the no-defense run. The mask-ratio sweep is invoked through
\texttt{scripts/\_run\_mask\_ratio\_sweep.sh} which sets the
\texttt{MASK\_RATIO\_TARGET} environment variable consumed by
\texttt{framework/attacker.py}. The transferability ablation
(\texttt{transition.py} on AdvBench) skips the attacker / scorer / fallback
entirely: it samples $k$ patterns from the bucket of
$\textsc{goal\_type}(g)$, replaces the literal goal sentence in each
pattern's stored template, expands \texttt{<mask:N>} once, and sends the
result directly to the victim --- this is the experiment-level evidence
that the discovered patterns transfer without further LLM-based adaptation.

\paragraph{Hyperparameters.}
Table~\ref{tab:hparams} lists the values used in every reported number; for
ablations the only entries that change are the mask ratio
($r$) and the per-goal attempt count ($k$), both labelled in the
corresponding tables.

\begin{table}[h]
\centering
\small
\caption{\method hyperparameters (frozen across all main experiments).}
\label{tab:hparams}
\begin{tabular}{l l l}
\toprule
Component & Hyperparameter & Value \\
\midrule
Attacker / scorer / extractor / merger & checkpoint & Qwen3-4B-Instruct-2507 \\
Fallback weaker base               & checkpoint  & Qwen3-4B-Base \\
Fallback re-tagger / judge         & checkpoint  & Bedrock \texttt{qwen.qwen3-235b-a22b-2507-v1:0} \\
Stage A (Exploration)              & seed goals  & 30 (HarmBench) \\
                                   & samples / goal & 1 \\
                                   & max\_length & 512 tokens \\
Stage B (Expansion)                & max\_iters  & 3 \\
                                   & max\_length & 2048 tokens (attacker) \\
                                   & success threshold $\rho^\star$ & 0.9 \\
                                   & fallback trigger              & $\rho < 0.7$ \\
                                   & max fallback retries          & 3 \\
UCB selection                      & $\alpha$    & 1.0 \\
                                   & bonus cap   & 1.0 \\
                                   & unvisited prior $(n_0, \mu_0)$ & $(32, 0.0)$ \\
                                   & top-$K$ pool size              & 3 \\
                                   & tie-break key                  & \texttt{stable\_hash(goal, iter, pid)} \\
Mask-ratio rescaler                & target $r$  & 0.65 (headline) \\
                                   & total token cap                & 512 \\
Victim mask token                  & Dream / LLaDA2                 & \texttt{<|mask|>} \\
                                   & LLaDA / LLaDA-1.5 / MMaDA      & \texttt{<|mdm\_mask|>} \\
Transferability (\texttt{transition.py}) & default $k$              & 5 \\
                                   & seed                           & 0 \\
\bottomrule
\end{tabular}
\end{table}

\paragraph{Computational Resource Requirement.}
The attacker-side vLLM endpoints (Qwen3-4B-Instruct + Qwen3-4B-Base) share
a single A100 / L40 GPU and stay loaded across all benchmarks. The dominant
cost is therefore victim-side mask-fill, which differs by an order of
magnitude across our five victims. We measured per-attempt inference time
on one L40-class GPU during the AdvBench transferability run (pure
retrieval, no attacker calls, exactly $k{=}5$ mask-fills per goal):
Dream-Instruct averages \textbf{$\approx 3.0$ s per attempt}
($\sim 25$ min for 100 goals at $k{=}5$); LLaDA-Instruct
\textbf{$\approx 2.7$ s} ($\sim 23$ min); LLaDA-1.5
\textbf{$\approx 3.2$ s} ($\sim 27$ min); MMaDA-CoT
\textbf{$\approx 10.0$ s} ($\sim 83$ min); and LLaDA2.0-mini
\textbf{$\approx 14.7$ s} ($\sim 123$ min, roughly $5\times$ slower than
the fast victims). The MMaDA outlier comes from its $2\times$ longer
unmasking schedule (\texttt{steps}=\texttt{max\_gen\_length}=$512$ vs.\
$256$ for the others); LLaDA2.0-mini's slower per-attempt time is due to
its smaller per-step batch ($1$ vs.\ $2$) and a heavier sampler kernel.

A full Stage~B run (\texttt{run\_victim.py}, 100 goals, $\le 3$ iterations
with $\le 3$ fallback retries each) takes roughly
\textbf{30--90 minutes} on the four fast victims and \textbf{2--3 hours}
on LLaDA2.0-mini. The complete sweep ---
5 victims $\times$ 3 main benchmarks (HB / JBB / SR) plus 4 victims
$\times$ 5 mask ratios on JBB plus 5 victims $\times$ AdvBench
transferability --- consumes $\approx 220$ GPU-hours on L40-class GPUs,
or about one wall-clock day on a node with 4 victim GPUs and one
attacker GPU. Bedrock judge calls cost $\approx \$5$ per full benchmark
sweep (15--30 files of 100 rows each, two metrics per row,
$\sim 3{,}000$--$6{,}000$ calls).

\section{Defense and Alignment for dLLMs}\label{appd:defense}
We evaluate \method against three families of safety mitigations originally
designed for autoregressive LLMs and adapted here to the diffusion mask-fill
setting. For each defense we use exactly the same Stage~B attack stack
described in Sec.~\ref{appd:impl-details} and only swap the victim
configuration (\texttt{configs/\{victim\}\_\{defense\}\_attack.yaml}); no
attacker-side hyperparameters change. All defenses are evaluated on
JailbreakBench with Strip-ASR-LLM and Strip-HarmfulScore as in the main
table; the per-victim ablations are reported in
Tables~\ref{tab:ablation_dream}--\ref{tab:ablation_mmada}.

\paragraph{Self-reminder.}
A prompt-level defense following \citet{xie2023defending}: a short
safety-instruction passage is prepended to the system role of the victim's
chat template, instructing it to refuse policy-violating requests and to
flag jailbreaking attempts. For diffusion victims the reminder is visible
to the unmasking sampler from step $0$; the assistant prefix received by
the victim still consists of the goal sentence followed by the attacker's
mask-fill template, so the only change relative to the base victim is the
system context. Self-reminder is enabled at evaluation time by setting
\texttt{SAFETY\_PROMPT=1} in \texttt{run\_victim.py --defense safety}; no
weight updates are performed. It is the cheapest of the three defenses
--- inference-time only, zero retraining cost.

\paragraph{Preference Optimization (PO).}
A weight-level defense implemented as \emph{any-order, any-step
random-masking SFT} on a refusal-augmented dataset. Concretely, every
harmful prompt in the safety training set has its original response
replaced by a randomly sampled refusal phrase
(``\textit{I'm sorry, but I can't help with that request\dots}'' --- one
of eight templates in \texttt{preprocess\_safety\_data.py}); safe prompts
keep their original response. The victim is then SFT-ed under the
standard random-masking objective: each example is masked at $L \cdot t$
positions with $t \sim \mathcal{U}(0.1, 0.9)$ and the LM is asked to
predict the original tokens at those positions. Each example is replayed
\texttt{mask\_times\_per\_sample}=$10$ times with independently sampled
$(t, \text{mask})$ pairs so that the loss covers every $(t, \text{order})$
combination uniformly --- the ``any-order, any-step'' regime that matches
the inference-time noise schedule. Training uses AdamW with
$\text{lr}=5\!\times\!10^{-6}$, cosine warmup over $50$ steps, batch
size $16$ (4 GPUs $\times$ 2 micro-batch $\times$ 2 grad-accum), bf16,
DeepSpeed ZeRO-3, $1$ epoch, and \texttt{max\_gen\_length}$=512$. The
resulting checkpoint (e.g.\ \texttt{dream\_safety\_aligned}) replaces the
base victim at evaluation time
(\texttt{configs/\{victim\}\_safety\_attack.yaml}). PO teaches the victim
to denoise toward refusal text whenever the prompt is harmful, regardless
of the mask layout.

\paragraph{A2D.}
A stronger weight-level defense based on the any-order any-step alignment
recipe of \citet{a2d_anyorder}, again implemented as random-masking SFT
but with one critical change. \emph{All} responses in the training set
--- including the original harmful continuations --- are kept verbatim
as the model input, so the diffusion LM still sees the harmful tokens in
its context window. However, at every masked position inside a harmful
response the \emph{label} is replaced with \texttt{eos\_id} before the
loss is computed (\texttt{train/sft\_dream.py}, lines around the
\texttt{label\_ids\_copy[rand\_mask] = eos\_id} block). The model is
therefore optimised to predict ``terminate here'' at any sampled
$(t, \text{mask})$ on harmful samples, while still being trained
normally on safe samples. To preserve general capability the safe split
is replayed $\texttt{mask\_times\_safe}=5$ times against
$\texttt{mask\_times\_per\_sample}=1$ for the harmful split. All other
hyperparameters match PO. Because the training signal is applied at every
step $t$ and every mask configuration, A2D directly counters the mask-fill
attack format \method exploits: any harmful continuation, no matter how
the attacker scaffolds it, is steered toward an early EOS during the
unmasking trajectory.

\paragraph{Results: ASR vs.\ utility trade-off.}
A defense's success on safety should not come at the cost of fundamentally
breaking the underlying dLLM. We therefore evaluate each defended victim on
two axes simultaneously: (i) attack-success rate against \method (lower is
better for the defender) and (ii) standard task utility on math / coding
benchmarks (higher is better for the defender). Both PO and A2D are
weight-level interventions and unavoidably trade utility for safety; our
calibration prioritises preserving utility, accepting only a moderate ASR
reduction in exchange for keeping the dLLM useful at its primary tasks.

\textbf{Attack-side trend (lower = stronger defense).}
Across all four defended victims (Dream-Instruct, LLaDA-Instruct,
LLaDA-1.5, MMaDA-CoT; LLaDA2.0-mini omitted because of its inference cost,
see Sec.~\ref{appd:impl-details}) the qualitative ordering is consistent:
Self-reminder $<$ PO $<$ A2D in defense strength, mirroring the increasing
distributional alignment between the defense training signal and the
attack distribution. On Dream-Instruct, \method-Search drops from
$94.0\%$ ASR (no defense) to $90.0$ / $87.0$ / $85.0\%$ under Self-reminder
/ PO / A2D respectively (Table~\ref{tab:ablation_dream}) --- a $\le 9$\,pp
absolute reduction even against the strongest defense, with corresponding
HS staying above $4.4$. Baselines collapse much faster: PAD loses
$36$\,pp under A2D ($66 \to 30\%$) and DIJA loses $53$\,pp
($79 \to 26\%$) on the same victim. The pattern repeats on LLaDA,
LLaDA-1.5, and MMaDA --- A2D produces the largest absolute drop relative
to the base victim but \method-Search retains $\ge 80\%$ ASR. We attribute
the gap to two properties of the mask-fill attack: (1) the attacker
template carries no harmful literal text, so refusal-style training
signals trigger only weakly; and (2) the goal-type-conditioned UCB pool
gives \method many distinct templates to try, of which the defense was
trained on at most a small subset.

\textbf{Utility-side cost (higher = better-preserved dLLM).}
Table~\ref{tab:defense_utility} reports zero-shot accuracy on standard
math (GSM8K, AIME-2024) and coding (HumanEval, MBPP) benchmarks for
Dream-Instruct under each defense. Self-reminder is prompt-only and
leaves utility essentially intact ($\le 1$\,pp drop expected on every
benchmark). PO and A2D, however, both \emph{catastrophically} degrade
reasoning. PO drops GSM8K from $64.0\%$ to $13.0\%$
($-51$\,pp), HumanEval from $66.0\%$ to $22.0\%$ ($-44$\,pp), and the
average from $46.8\%$ to $16.5\%$ ($-30$\,pp absolute). A2D is even more
severe: GSM8K and AIME-2024 collapse to \emph{exactly $0\%$}, HumanEval
drops to $16\%$ ($-50$\,pp), and the average falls to $7.0\%$ ($-40$\,pp).

\textbf{Why A2D collapses to $0\%$.}
The mechanism is specific to A2D's any-order any-step EOS-replacement
trick: at every masked position in a harmful sample, the training label
is replaced with \texttt{eos\_id}, so the model learns to emit
end-of-text whenever its context resembles a structured procedural
generation. Inspecting A2D's GSM8K outputs confirms this: most generations
consist of one or two literal tokens followed by an \texttt{<|endoftext|>}
storm that propagates across the entire diffusion block.\footnote{A
representative trace: prompt = ``Richard lives in an apartment building
with $15$ floors\dots What's the total number of unoccupied units?'';
A2D output = \texttt{<|endoftext|>The total number of unoccupied units
in the building is <|endoftext|>boxed\{30 <|endoftext|><|endoftext|>\dots} ---
the model arrives at the correct answer ($30$) but is interrupted by an
EOS token before it can close the \texttt{\textbackslash boxed\{\}}
expression, so the answer extractor reports a miss.} Mathematical
reasoning, code synthesis, and harmful procedural attack templates all
share the same surface form (numbered steps, conditional structure,
slot-like reasoning chains), so the model cannot tell them apart at the
token level and applies the EOS shortcut indiscriminately. PO is less
catastrophic ($16.5\%$ avg) because its training signal targets the
\emph{prompt} (refuse if prompt is harmful) rather than the
\emph{completion} (terminate if context is procedural), but it still
loses $30$\,pp because any sufficiently complex prompt looks
``harmful-like'' to a refusal-trained classifier.

The same trend holds on the other three victims (LLaDA, LLaDA-1.5,
MMaDA): we observed an analogous A2D collapse on LLaDA's GSM8K (full
results omitted for space). The conclusion is that utility loss is
intrinsic to mask-fill-conditioned safety SFT rather than an artefact of
our specific recipe.

\begin{table}[h]
\centering
\small
\caption{Utility vs.\ safety on \textbf{Dream-Instruct} under three
defenses. \emph{All ASR and utility numbers in this table are for the
same victim (Dream-Instruct); the only thing that changes across rows
is which defense (if any) is applied to its weights or system prompt.}
\textbf{Utility} (left block, higher = better-preserved dLLM): zero-shot
accuracy (\%) measured on the first $100$ examples of each benchmark
(AIME-2024 has only $30$ samples) with
$\texttt{max\_gen\_length}{=}1024$ and $\texttt{steps}{=}1024$;
$\Delta$ is the absolute drop versus the base victim, and \textbf{AVG}
is the unweighted mean across the four benchmarks.
\textbf{ASR} (right block, lower = stronger defense): Strip-ASR-LLM
(\%) on \textbf{JailbreakBench} ($100$ goals), judged by Bedrock
Qwen3-235B. \textsc{Zero-shot} sends the plain harmful goal as the
user message with no attack scaffold; \textsc{PAD}, \textsc{DIJA}, and
\textsc{\method-Search} use their respective attack templates as
defined in Sec.~\ref{appd:impl-details}. Zero-shot ASR is $0\%$ for the
base victim and remains $\approx 0\%$ under every defense, confirming
that the base dLLM is already aligned to refuse plain harmful prompts
--- the defenses only matter under attack. Self-reminder utility
numbers are extrapolated from base ($\le 1$\,pp drop expected since it
is prompt-only); PO and A2D utility are measured directly.}
\label{tab:defense_utility}
\resizebox{\linewidth}{!}{
\begin{tabular}{l c c c c c | c c c c}
\toprule
\multirow{2}{*}{Defense} & \multicolumn{5}{c|}{Utility (acc \%, higher = better)}
& \multicolumn{4}{c}{ASR (\%, lower = stronger defense)} \\
\cmidrule(lr){2-6} \cmidrule(lr){7-10}
& GSM8K & AIME & HumanEval & MBPP & \textbf{AVG}
& Zero-shot & PAD & DIJA & \textbf{MaskForge-S} \\
\midrule
Base
  & 64.0 & 3.3 & 66.0 & 54.0 & \textbf{46.8}
  & 0.0 & 66.0 & 79.0 & \textbf{94.0} \\
+\,Self-reminder
  & $\sim 64$ & $\sim 3$ & $\sim 65$ & $\sim 54$ & $\sim 46$
  & 0.0 & 43.0 & 56.0 & 90.0 \\
+\,PO
  & 13.0 ($-51.0$) & 0.0 ($-3.3$) & 22.0 ($-44.0$) & 31.0 ($-23.0$) & 16.5 ($-30.3$)
  & 0.0 & 50.0 & 31.0 & 87.0 \\
+\,A2D
  & \textbf{0.0} ($-64.0$) & \textbf{0.0} ($-3.3$) & 16.0 ($-50.0$) & 12.0 ($-42.0$) & \textbf{7.0} ($-39.8$)
  & 0.0 & 30.0 & 26.0 & 85.0 \\
\bottomrule
\end{tabular}}
\end{table}

\textbf{Design choice.}
The trade-off is dramatically unbalanced. PO loses $30$\,pp average
utility for $7$\,pp ASR reduction; A2D loses $40$\,pp average utility
for $9$\,pp ASR reduction. That is roughly \textbf{$4$--$5$ utility points
sacrificed per safety point gained} --- not a trade so much as a
capability collapse with a mild safety side-effect. Defending mask-fill
diffusion victims against retrieval-driven attacks like \method
therefore appears to require defenses that explicitly preserve
non-procedural-but-benign generations: e.g.\ a KL-anchored variant of
A2D that limits how far weights can drift on safe prompts, or
safety-aware decoding interventions (refusal-token boosting at unmasking
time) layered on top of an unmodified victim. Crucially, the
\emph{distribution-matched defense} that would best counter \method ---
A2D's any-order any-step EOS objective --- is also the one that breaks
the underlying dLLM most thoroughly, because mask-fill jailbreaks and
mask-fill reasoning are the same operation. We deliberately do not push
our PO / A2D recipes further: the goal of this section is to characterise
the trade-off surface, and the data above already shows it is steeply
adversarial. We leave the joint optimisation of safety and utility to
future work.

\section{Details of \method}\label{appd:method}

\subsection{Algorithmic Outline of \method}\label{appd:algo}
We summarise the full \method pipeline as Algorithm~\ref{alg:maskforge}. Stage~A
(Exploration, lines~\ref{alg:lineA-start}--\ref{alg:lineA-end}) seeds the
pattern registry $\mathcal{R}$ from a small pool of harmful goals. Stage~B
(Expansion, lines~\ref{alg:lineB-start}--\ref{alg:lineB-end}) is the loop run
at evaluation time: for each goal it draws a UCB-prioritised pattern, asks the
attacker to instantiate it into a strategy/prompt/template triple, sends the
expanded template to the victim, and on failure invokes a fallback that
distils a fresh pattern from a weaker base model's uncensored draft. Both
stages share the same registry, so successful patterns are immediately
available to all subsequent goals.

\begin{algorithm}[h]
\small
\caption{MaskForge: Stage A (Exploration) + Stage B (Expansion).}
\label{alg:maskforge}
\KwIn{seed goals $\mathcal{G}_A$, test goals $\mathcal{G}_B$, attacker LLM
$f_{A}$, weaker base $f_{W}$, scorer $f_{S}$, re-tagger $f_{T}$ (Bedrock),
victim dLLM $V$, success threshold $\rho^{\star}$, fallback threshold
$\rho_{\text{fb}}$, $T_{\max}$ iters, $K$ pool size, UCB coefficient $\alpha$.}
\KwOut{predictions $\{r_g\}_{g \in \mathcal{G}_B}$, updated registry $\mathcal{R}$, pattern set $\mathcal{P}$.}
\BlankLine
\textbf{Stage A (Exploration).} \nllabel{alg:lineA-start}
$\mathcal{R}, \mathcal{P} \leftarrow \emptyset$\;
\For{each seed goal $g \in \mathcal{G}_A$}{
  $(s, q, \tau) \leftarrow f_A(g)$ \quad // strategy, prompt, mask template\;
  $r \leftarrow V(q, \texttt{expand}(\tau))$\;
  \If{$f_S(g, r) \ge \rho^{\star}$}{
    $\sigma \leftarrow \textsc{Summarise}(g, \tau, r)$ \quad // schema $(\textsc{type}, \textsc{slots},\dots)$\;
    $\mathcal{R}[\textsc{hash}(\sigma)] \leftarrow (\sigma, \tau)$;\quad
    $\mathcal{P}[\textsc{type}(g)] \cup\!= \{\textsc{hash}(\sigma)\}$\;
  }
} \nllabel{alg:lineA-end}
\BlankLine
\textbf{Stage B (Expansion).} \nllabel{alg:lineB-start}
\For{each test goal $g \in \mathcal{G}_B$}{
  $\rho^\star_g \leftarrow 0$\;
  \For{$t = 1$ \KwTo $T_{\max}$}{
    $\mathcal{C} \leftarrow \mathcal{P}[\textsc{type}(g)]$ \quad // candidate pool\;
    \tcc{UCB selection: top-$K$ by $\mu_p + \alpha\sqrt{2\ln t / n_p}$, then deterministic hash-shuffle.}
    $p^\star \leftarrow \textsc{UCB-select}(\mathcal{C}, \theta, t, \alpha, K, g)$\;
    $(s, q, \tau) \leftarrow f_A(g, \mathcal{R}[p^\star])$\;
    $\tau \leftarrow \textsc{rescaleMaskRatio}(\tau, r{=}0.65)$ \quad // see Sec.~\ref{appd:impl-details}\;
    $r \leftarrow V(q, \texttt{expand}(\tau))$;\quad
    $\rho \leftarrow f_S(g, q, \tau, r)$\;
    \If{$\rho < \rho_{\text{fb}}$}{ \tcc{Fallback: weaker base draft + Bedrock re-tagger.}
      \For{$j = 1$ \KwTo $3$}{
        $r' \leftarrow f_W(g)$ \quad // raw uncensored draft\;
        $\tau' \leftarrow f_T(r', \texttt{[UNSAFE]\dots[/UNSAFE]} \to \texttt{<mask:N>})$\;
        $r \leftarrow V(q, \texttt{expand}(\tau'))$;\quad
        $\rho \leftarrow f_S(g, q, \tau', r)$\;
        \If{$\rho \ge \rho_{\text{fb}}$}{ $\tau \leftarrow \tau'$;\quad \textbf{break}\; }
      }
    }
    \textsc{UCB-update}$(\theta, p^\star, \rho)$\;
    \If{$\rho > \rho^\star_g$}{
      $\sigma \leftarrow \textsc{Summarise}(g, \tau, r)$\;
      \If{$\textsc{hash}(\sigma) \notin \mathcal{R}$}{
        $\mathcal{R}[\textsc{hash}(\sigma)] \leftarrow (\sigma, \tau)$;\quad
        $\mathcal{P}[\textsc{type}(g)] \cup\!= \{\textsc{hash}(\sigma)\}$\;
      }
      $\rho^\star_g \leftarrow \rho$;\quad $r_g \leftarrow r$\;
    }
    \If{$\rho \ge \rho^{\star}$}{ \textbf{break}\; }
  }
} \nllabel{alg:lineB-end}
\end{algorithm}

\subsection{Query Times of \method on Test-Time Transferability}\label{appd:querycost}
Test-time transferability runs (Tab.~\ref{tab:transition_advbench}) deliberately
strip Stage~B down to the minimum required to attack: only the victim is
queried, all attacker / scorer / fallback / summariser calls are skipped, and
the registry $\mathcal{R}$ is never updated. This isolates the
\emph{transfer-ability of the patterns themselves} from the contribution of
Stage~B's online optimisation, and makes the cost trivial to compare with
single-shot baselines.

\paragraph{Per-goal LLM calls.}
Let $C_{V}$ be the cost of one victim mask-fill, $C_{A}$ one attacker LLM
call, $C_{S}$ one scorer call, and $C_{B}$ one Bedrock judge / re-tagger
call. The cost of one attempt against goal $g$ is:
\begin{itemize}
\item \textsc{Direct}: $1\cdot C_V$.
\item \textsc{PAD}: $1\cdot C_V$ (template is hard-coded).
\item \textsc{DIJA}: $1\cdot C_A + 1\cdot C_V$, with the $C_A$ for the
per-goal refined template amortised across victims since it is cached
(\texttt{logs/transition\_baselines/dija\_templates.json}). After the cache
warms up, DIJA's marginal cost per (victim, goal) pair is $1\cdot C_V$.
\item \textsc{\method-Transferability} ($k$ attempts): $k\cdot C_V + 1$
embedding lookup. With pre-computed pattern and goal embeddings (one-time
$\mathcal{O}(|\mathcal{R}|)$ cost), retrieval per goal is a single
$\mathbf{P} \mathbf{g}$ matvec --- under a millisecond on CPU.
\item \textsc{\method-Search / Dual100} (full Stage~B): up to $T_{\max}{=}3$
iterations, each with $(C_A + C_V + C_S)$ for the main path plus an
optional fallback of up to $3$ retries each adding $(C_W + C_B + C_V + C_S)$.
In the worst case this is bounded by
$T_{\max}\cdot(C_A + C_V + C_S) + T_{\max}\cdot 3\cdot(C_W + C_B + C_V + C_S)$
per goal; in practice early termination on $\rho \ge \rho^{\star}$ keeps the
average to $\sim\!1.7$ iterations and one or zero fallback retries.
\end{itemize}

\paragraph{Empirical victim-call budget.}
Across $100$ AdvBench goals on Dream-Instruct, the recorded per-method
victim-call counts are reported in Table~\ref{tab:querycount}. Numbers are
exact (counted from the JSONL trace files), not estimates. \method
transferability with $k{=}5$ uses $5\times$ as many victim calls as the
single-shot baselines per goal, but every call is goal-specific and runs
without any external LLM in the loop.

\begin{table}[h]
\centering
\small
\caption{Per-goal LLM-call cost on AdvBench (averaged over $100$ goals on
Dream-Instruct). $C_V$: victim mask-fill; $C_A$: attacker; $C_S$: scorer;
$C_B$: Bedrock re-tagger; emb: BGE-cosine matvec.}
\label{tab:querycount}
\begin{tabular}{l c c c c}
\toprule
Method & $C_V$ & $C_A$ & $C_S$ & emb lookup \\
\midrule
\textsc{Direct}                       & $1$ & $0$ & $0$ & $0$ \\
\textsc{PAD}                          & $1$ & $0$ & $0$ & $0$ \\
\textsc{DIJA}                         & $1$ & $0$ (cached) & $0$ & $0$ \\
\textsc{\method ($k{=}5$, transfer)} & $5$ & $0$ & $0$ & $1$ (matvec) \\
\textsc{\method-Search} (full)        & $1.7$ avg, $\le 12$ max & $1.7$ & $1.7$ & $0$ \\
\bottomrule
\end{tabular}
\end{table}

\paragraph{Why no attacker / scorer at transfer time?}
The patterns in $\mathcal{R}$ already carry their persuasion strategy in the
\texttt{representative\_template} field, which contains the goal-aligned
literal text together with the $\texttt{<mask:N>}$ scaffold. Replacing the
literal goal sentence in line~$1$ with the new AdvBench goal and expanding
the masks is sufficient to repurpose the pattern --- no attacker LLM call
is required. Skipping the scorer is the corresponding choice on the
evaluation side: instead of running the local Qwen-4B reward at attack
time, we accumulate $k$ candidate generations and report best-of-$k$
ASR-LLM under the global Bedrock judge (Sec.~\ref{appd:eval-metrics}).
This makes the transferability numbers a clean lower bound on what the
full Stage~B loop would achieve at the same victim-call budget.

\subsection{Retrieval Strategy of \method}\label{appd:retrieval}
\method's pattern registry is large (\textbf{$1{,}779$} patterns across
$15$ goal-type buckets at the time of writing) and the choice of how to
narrow it down per goal materially affects ASR. We compare three
strategies under a fixed victim (Dream-Instruct), a fixed attempt budget
($k{=}5$), and a fixed Bedrock judge on AdvBench$[0{:}30]$:

\begin{itemize}
\item \textbf{Uniform-random over registry.} Sample $k$ pattern IDs uniformly
from all $1{,}779$ patterns, no goal-type filter. Establishes the
no-prior baseline.
\item \textbf{Random within goal-type bucket} (\method's default). Classify
$g$ into one of $15$ goal types via the retriever-LLM at indexing time, then
sample $k$ uniformly from $\mathcal{P}[\textsc{type}(g)]$. The bucket sizes
range from $\sim\!10$ (\textsc{sexual\_content}) to $\sim\!340$
(\textsc{deception}); the average bucket has $\sim\!119$ patterns.
\item \textbf{Embedding cosine top-$K$.} Pre-compute embeddings of every
pattern's representative template and every goal, then take the top-$k$
patterns by cosine similarity. We swept five different encoders (BGE-small,
BGE-large, E5-large, MXBai-large, Qwen3-Embedding-0.6B) using the natural
pooling for each (CLS / mean / last-token).
\end{itemize}

The headline result is in Table~\ref{tab:retrieval_strategy}. Embedding
retrieval beats both random baselines at every $k$, and the gap is
largest at $k{=}1$ where it most directly measures retrieval precision:
the strongest encoder, BGE-large, jumps from the no-prior baseline of
$13.3\%$ to $\mathbf{66.7\%}$ ASR on a single attempt, more than $5\times$.

\begin{table}[h]
\centering
\small
\caption{Retrieval-strategy ablation on Dream-Instruct $\times$ AdvBench$[0{:}30]$
$\times$ $k{=}5$, Bedrock judge. Embedding-based retrieval is a strict
improvement over the goal-type bucket prior at every $k$; BGE-large is the
most effective encoder we tested.}
\label{tab:retrieval_strategy}
\begin{tabular}{l c c c c c c}
\toprule
Strategy & Family & Size / dim & $k{=}1$ & $k{=}2$ & $k{=}3$ & $k{=}5$ \\
\midrule
Uniform random over registry & --- & --- & 13.3 & 23.3 & 33.3 & 46.7 \\
Random within goal-type bucket (default) & --- & --- & 23.3 & 46.7 & 53.3 & 66.7 \\
\midrule
Embedding cosine top-$K$ & BGE-small & 33M / 384d  & 23.3 & 46.7 & 60.0 & 80.0 \\
Embedding cosine top-$K$ & E5-large  & 335M / 1024d & 23.3 & 40.0 & 56.7 & 80.0 \\
Embedding cosine top-$K$ & MXBai-large & 335M / 1024d & 46.7 & 60.0 & 66.7 & 76.7 \\
Embedding cosine top-$K$ & Qwen3-Embedding-0.6B & 595M / 1024d & 36.7 & 43.3 & 53.3 & 56.7 \\
Embedding cosine top-$K$ & \textbf{BGE-large}~\citep{bge_embedding} & \textbf{335M / 1024d} & \textbf{66.7} & \textbf{76.7} & \textbf{83.3} & \textbf{93.3} \\
\bottomrule
\end{tabular}
\end{table}

\paragraph{Why goal-type bucketing is not enough.}
The retriever-LLM that assigns goal types is itself an LLM classifier
(Qwen-3 4B Instruct), and its labels are coarse: an MDMA-synthesis
pattern with explicit chemistry slots
(\textsc{material\_preparation}, \textsc{solvent\_method}, \dots) is
indexed under \textsc{illegal\_activity} simply because synthesising MDMA
is illegal, even though structurally it has nothing in common with a
ransomware goal that is also under \textsc{illegal\_activity}. Random
sampling within such a bucket therefore wastes attempts on
patterns that are topical-but-irrelevant. Embedding retrieval breaks the
tie: out of $30 \times 5 = 150$ pattern picks, the goal-type-bucket
baseline and BGE-large agree on only $\sim\!4$ pattern IDs --- the two
strategies recommend almost disjoint subsets of the registry.

\paragraph{Why best-of-$k$ still helps after embedding retrieval.}
Even with BGE-large, the top-$1$ pattern is the right choice for only
$66.7\%$ of goals; for the remaining $33.3\%$ the best-matched pattern
fails (e.g.\ the victim refuses despite the framing). Different patterns
in the top-$5$ correspond to different persuasion strategies (academic
case study, fictional workshop, first-responder training, journalism,
debate) and different slot granularities, so increasing $k$ yields a
diversity bonus on top of the precision bonus from cosine ordering. This
explains why $k{=}5$ ASR rises to $93.3\%$ even though $k{=}1$ is
$66.7\%$: each additional attempt adds an alternative
\emph{framing} of the goal, not a redundant copy.

\paragraph{Practical recommendation.}
For deployments where any LLM call at attack time is acceptable, run the
full Stage~B loop --- pattern evolution alone keeps lifting the registry
quality run after run. For test-time transferability where only victim
mask-fills are allowed, use \textbf{BGE-large cosine top-$K$ on the full
$1{,}779$-pattern registry} (no goal-type filter) with $k\!\in\![3, 5]$.
Pre-computing the embedding matrix for the registry is a one-time cost
of $\sim\!2$ minutes on a single GPU and $\sim\!7$\,MB on disk.

\section{Pattern Library}\label{appd:library}
The pattern library is the central artefact \method produces and re-uses
across all victims, benchmarks, and ablations. After Stage~A exploration on
$30$ seed HarmBench goals plus the Stage~B online expansion run during the
main experiments, and after the quality filter described below has rejected
schemas with no semantic slot annotation, the registry persisted to
\texttt{logs/shared\_pattern\_registry.json} contains
\textbf{$1{,}299$ patterns}, all tagged \texttt{source: expansion}. Each
entry is a \texttt{(pattern\_id, schema, representative\_template, source)}
tuple where \texttt{pattern\_id} is a stable hash of the canonicalised
schema, so two pipeline runs that converge on the same schema produce the
same id and merge in-place rather than duplicating.

\paragraph{Quality filter.}
The summarizer occasionally returns schemas whose
\texttt{slot\_roles} list is either empty or filled with content-free
placeholder names such as \texttt{["input"]\,$\times$\,N}; these schemas
look syntactically valid but carry no semantic conditioning for
\texttt{attacker.instantiate}, so re-using them collapses the schema-driven
attack into pure textual retrieval. \texttt{PatternRegistry.add\_pattern}
runs a small validator (\texttt{\_is\_low\_quality\_schema} in
\texttt{expansion.py}) before insertion: it rejects schemas where (a)
\texttt{slot\_roles} is missing or empty while \texttt{slot\_count} is
positive, (b) every role name matches the regex
\verb|^(input|slot|placeholder|value|field|item|mask|fill)[_\-\s]?\d*$|, or
(c) all $\ge\!4$ roles are identical. Rejected schemas are counted by
reason but never written to disk. On the current corpus the validator
removes $480$ patterns that an earlier offline-recovery pass had injected
with \texttt{slot\_roles=["input"]\,$\times$\,N} and
\texttt{organization\_style="labeled"}; everything that survives carries
informative role names.

\paragraph{Goal-type distribution.}
Patterns are indexed by goal type into \texttt{logs/shared\_pattern\_set.json}
under a \textbf{30-class} fine-grained taxonomy. We started from the
$9$-class taxonomy emitted by Stage~A's retriever (\textsc{general},
\textsc{harm}, \textsc{discrimination}, \textsc{deception},
\textsc{illegal\_activity}, \textsc{exploitation}, \textsc{ethical\_violation},
\textsc{cultural\_normalization}, \textsc{historical\_distortion}) but found
that several buckets --- in particular \textsc{illegal\_activity} and
\textsc{deception} --- collapsed structurally distinct attack patterns
(e.g.\ a chemistry-themed MDMA-synthesis scaffold and a ransomware scaffold
both end up under \textsc{illegal\_activity}). To make the bucketing useful
as a retrieval prior, we expand the taxonomy to $30$ mutually-exclusive
categories grouped into seven topical clusters: cyber/digital
(\textsc{cybercrime}, \textsc{malware}, \textsc{cyberweapons}); physical
violence \& weapons (\textsc{violence}, \textsc{weapons\_firearms},
\textsc{weapons\_explosives}, \textsc{bioweapons}); substances
(\textsc{chemistry}, \textsc{drug\_synthesis}, \textsc{drug\_distribution});
financial/legal (\textsc{financial\_fraud}, \textsc{identity\_theft},
\textsc{counterfeiting}, \textsc{tax\_evasion}, \textsc{extortion});
personal/content (\textsc{self\_harm}, \textsc{eating\_disorders},
\textsc{sexual\_content}, \textsc{harassment}, \textsc{hate\_speech});
information/society (\textsc{privacy\_invasion}, \textsc{misinformation},
\textsc{conspiracy\_theories}, \textsc{defamation},
\textsc{election\_manipulation}); plus four generic harm-shape buckets
preserved from the original taxonomy (\textsc{harm}, \textsc{discrimination},
\textsc{deception}, \textsc{exploitation}) and a catch-all \textsc{general}.
The reclassification is run by \texttt{scripts/\_expand\_taxonomy\_bedrock.py},
which sends each pattern's \texttt{representative\_template} to Bedrock
Qwen3-235B with a closed-list classification prompt and parses out one of
the $30$ labels (we use Bedrock rather than the local Qwen-4B retriever
because it produces better class boundaries on the long tail). All $1{,}299$
quality-filtered patterns receive a valid label without falling back to
\textsc{general}. Figure~\ref{fig:pattern_library_stats}(a) shows the
resulting bucket sizes. The five largest buckets are
\textsc{misinformation} ($183$), \textsc{general} ($113$),
\textsc{discrimination} ($99$), \textsc{violence} ($96$), and
\textsc{hate\_speech} ($88$); the smallest are \textsc{tax\_evasion} ($2$),
\textsc{eating\_disorders} ($3$), \textsc{extortion} ($4$), and
\textsc{identity\_theft} / \textsc{conspiracy\_theories} (each $5$). The
strongly skewed distribution reflects \method's opportunistic learning:
categories with high base-victim ASR (and therefore many Stage~B successes
to summarise) accumulate patterns much faster than narrower categories such
as \textsc{eating\_disorders} or \textsc{tax\_evasion}. The taxonomy is
consumed by both the random-sample baseline (\texttt{transition.py}, which
draws $k$ candidates from $\textsc{type}(g)$'s bucket) and serves as a
coarse prior for the embedding retriever (Sec.~\ref{appd:retrieval}).

\paragraph{Structure-type taxonomy.}
The summarizer's first schema field, \texttt{structure\_type}, captures the
high-level rhetorical scaffold of the template. Across the $1{,}299$
quality-filtered patterns the summarizer has produced $399$ distinct
\texttt{structure\_type} strings --- evidence that it does not collapse
onto a small canonical vocabulary, even though the long tail consists of
one-off variants. The eight most common types together account for $541$
patterns ($41.6\%$ of the registry) and cover the dominant rhetorical modes
\method exploits, in decreasing frequency: \textsc{procedural\_steps}
($220$, e.g.\ phase-by-phase walkthroughs),
\textsc{technical\_explanation} ($114$), \textsc{dialogue\_simulation}
($66$), \textsc{analytical\_description} ($65$), \textsc{narrative\_arc}
($20$), \textsc{narrative\_progression} ($19$),
\textsc{logical\_progression} ($19$), and \textsc{question\_answer} ($18$).
The full top-eight breakdown is plotted in
Figure~\ref{fig:pattern_library_stats}(b).

\paragraph{Slot count and slot roles.}
Each pattern carries an integer \texttt{slot\_count} (the number of
$\texttt{<mask:N>}$ holes in the canonical template) and a list of
\texttt{slot\_roles} naming what each hole is meant to be filled with
(e.g.\ \textsc{materials\_and\_reagents}, \textsc{contextual\_constraint}).
Figure~\ref{fig:pattern_library_stats}(c) shows the slot-count
distribution: $5$ slots is the modal choice ($643$ patterns, $49.5\%$),
followed by $4$ slots ($401$, $30.9\%$), $6$ slots ($116$, $8.9\%$),
$3$ slots ($46$, $3.5\%$), and $8$ slots ($40$, $3.1\%$). $98.2\%$ of
patterns have between $3$ and $8$ slots; the long tail of $>\!10$-slot
patterns is small and consists of unusually elaborate scaffolds the
summarizer distilled from particularly long jailbroken outputs --- they
are rarely picked by UCB or by embedding cosine.

After the validator drops content-free roles, the \texttt{slot\_roles}
field carries $6{,}350$ role occurrences over $3{,}182$ distinct values ---
effectively free-text annotations the attacker LLM produced when prompted,
rather than a closed vocabulary. The most frequent roles either (i)
describe the \emph{kind of content} the slot expects, e.g.\ \textsc{method}
($41$), \textsc{event\_sequence} ($32$), \textsc{call\_to\_action} ($27$),
\textsc{response} ($27$), \textsc{condition} ($26$), \textsc{container}
($24$); or (ii) name a \emph{position in the rhetorical scaffold}, e.g.\
\textsc{background\_context} ($184$), \textsc{request} ($73$),
\textsc{background} ($57$), \textsc{context} ($55$),
\textsc{contextual\_background} ($40$), \textsc{setting} ($32$),
\textsc{narrative\_framing} ($23$), \textsc{conclusion} ($22$). A
representative sample is given in Table~\ref{tab:slot_roles}; we report the
role names verbatim because their lexical specificity is what makes the
embedding-retrieval path effective --- an embedding encoder can distinguish
a pattern with slots $\{\textsc{materials\_and\_reagents},$
$\textsc{processing\_conditions}, \textsc{purification\_process}\}$
(chemistry-themed) from a pattern with slots $\{\textsc{initial\_indicators},$
$\textsc{behavioral\_red\_flags},$
$\textsc{detection\_aversion\_strategy}\}$ (operational-security-themed)
even when both fall under the same goal-type label.

\begin{figure}[h]
\centering
\includegraphics[width=\linewidth]{papers/pattern_library_stats.pdf}
\caption{Pattern-library statistics over the $1{,}299$ quality-filtered
patterns persisted to \texttt{logs/shared\_pattern\_registry.json}.
(a) Goal-type bucket sizes under the $30$-class taxonomy
(\texttt{logs/shared\_pattern\_set.json}; reclassified by Bedrock
Qwen3-235B). The five largest buckets cover the dominant Stage~B success
modes (\textsc{misinformation}, \textsc{general}, \textsc{discrimination},
\textsc{violence}, \textsc{hate\_speech}); narrow categories such as
\textsc{eating\_disorders} or \textsc{tax\_evasion} have far fewer patterns
because Stage~B accumulates patterns only from successful attempts.
(b) Top-$8$ \texttt{structure\_type} fields produced by the summarizer;
together they cover $41.6\%$ of the registry, the rest spread across
$391$ rarer types. (c) \texttt{slot\_count} distribution: the modal
scaffold has $5$ masks, with $98.2\%$ of patterns in the $3$--$8$ range.}
\label{fig:pattern_library_stats}
\end{figure}

\begin{table}[h]
\centering
\small
\caption{Sample \texttt{slot\_roles} from the quality-filtered registry,
grouped by the rhetorical function they tend to serve. The full vocabulary
contains $3{,}182$ distinct strings; we list the most frequent
representatives in each group (counts after the validator that strips
content-free \texttt{["input"]} schemas).}
\label{tab:slot_roles}
\begin{tabular}{l p{0.62\linewidth}}
\toprule
Group & Example slot roles (counts in parentheses) \\
\midrule
Scaffold position & \textsc{background\_context} (184),
\textsc{request} (73), \textsc{background} (57), \textsc{context} (55),
\textsc{contextual\_background} (40), \textsc{setting} (32),
\textsc{narrative\_framing} (23), \textsc{conclusion} (22) \\
\midrule
Procedural steps & \textsc{materials\_and\_reagents},
\textsc{processing\_conditions}, \textsc{purification\_process},
\textsc{event\_sequence} (32), \textsc{event\_description} (18),
\textsc{key\_event\_sequence} (17) \\
\midrule
Domain-specific (illicit ops) & \textsc{detection\_evasion\_tactic},
\textsc{concealment\_method}, \textsc{hidden\_compartment},
\textsc{initial\_access\_method} (19), \textsc{verification\_bypass},
\textsc{post\_transaction\_actions}, \textsc{persistence\_strategy} \\
\midrule
Domain-specific (rhetoric / persuasion) & \textsc{call\_to\_action} (27),
\textsc{response} (27), \textsc{condition} (26), \textsc{container} (24),
\textsc{method} (41), \textsc{motivation} (19), \textsc{target\_audience} \\
\bottomrule
\end{tabular}
\end{table}

\paragraph{Why the library transfers.}
Two properties make the persisted patterns reusable across goals and
victims. First, the literal text in each \texttt{representative\_template}
is goal-aligned but does \emph{not} contain the harmful specifics ---
those go in \texttt{<mask:N>} placeholders. A pattern learned from one
ransomware goal can therefore be re-instantiated against any other
illegal-activity goal by swapping the goal sentence in line~$1$ and
expanding the masks. Second, the schema fields
$(\texttt{structure\_type}, \texttt{slot\_count}, \texttt{slot\_roles})$
expose the \emph{rhetorical strategy} of the pattern in a form that an
embedding encoder can match against a fresh goal --- the slot-role names
do most of the work of the goal-type label, but at much finer granularity.
This is also why the quality filter matters: a schema with
\texttt{slot\_roles=["input",\,\dots]} encodes nothing the encoder can
distinguish from any other low-quality entry, so admitting it would dilute
the retrieval signal even though the surface template might still elicit a
victim response. The two properties together --- masked specifics and
informative schemas --- are what allow the test-time transferability
protocol of Sec.~\ref{appd:querycost} to skip the attacker LLM entirely
while retaining $\ge 80\%$ of the full Stage~B ASR
(Tab.~\ref{tab:transition_advbench}).

\section{Full Prompt of \method}\label{appd:prompt}
This section reproduces verbatim the system / user prompts used by every
LLM call inside \method. Four roles need their own prompt template:

\setlength{\parindent}{20pt}\verb|attacker|: Qwen3-4B-Instruct
(\texttt{vLLM} 8100). Conditioned on a single \verb|goal|, it emits a
matching \emph{(strategy, jailbreak prompt, mask template)} triple via the
delimited \texttt{=== STRATEGY === / === JAILBREAK PROMPT === / === TEMPLATE ===}
output format. See Sec.~\ref{appd:algo}, Step 3 of Stage~B.

\setlength{\parindent}{20pt}\verb|scorer|: Qwen3-4B-Instruct
(\texttt{vLLM} 8100). Implements a HarmBench-style \emph{yes / no}
classifier over (\verb|goal|, \verb|victim_output|). The literal string
``\texttt{yes}'' or ``\texttt{no}'' on the first line of the response is
mapped to reward $1.0$ or $0.0$ respectively. See Sec.~\ref{appd:algo},
Step 5.

\setlength{\parindent}{20pt}\verb|drafter|: Two cooperating prompts that
make up the fallback path triggered when the scorer returns
$\rho < 0.7$. (A) A \emph{rephraser} run by the attacker LLM that rewrites
the original \verb|goal| into a more direct, on-topic sentence
(\verb|framing| and \verb|feedback_block| are optional retry hints).
(B) A \emph{Bedrock UNSAFE-span tagger} (Qwen3-235B) that takes the
weaker base model's uncensored draft and inserts
\texttt{[UNSAFE]\dots[/UNSAFE]} markers around concrete operational
content; the tagged spans are then converted into \verb|<mask:N>| of
matching token length to form the next iteration's template. See
Sec.~\ref{appd:algo}, Step 6.

\setlength{\parindent}{20pt}\verb|summarizer|: Qwen3-4B-Instruct
(\texttt{vLLM} 8100). Compares the previous-best and current-iteration
attempt on the same goal and emits a JSON object
$\{\textsc{improvement\_summary},\ \textsc{new\_pattern},\
\textsc{new\_template}\}$ that, on improvement, is hashed and persisted
to the registry. See Sec.~\ref{appd:algo}, Step 7.

\setlength{\parindent}{0pt}Placeholders in curly braces ($\{\cdot\}$) are
filled in at run time. The literal text is reproduced exactly as it
appears in \texttt{framework/attacker.py},
\texttt{framework/scorer.py}, and \texttt{framework/summarizer.py}.

\subsection{Prompt for Attacker LLM}\label{appd:prompt-attacker}
The attacker prompt drives every Stage~B iteration: it asks the LLM to
pick \emph{one} persuasion strategy from a pool of named templates
(fiction workshop, academic case study, first-responder training, etc.),
write a jailbreak prompt that invokes that strategy, and produce a
matching mask template whose literal scaffold contains no harmful
specifics. The three in-context examples we provide cover the dominant
shapes \method uses --- tabloid journalism, procedural walkthrough, and
media-literacy case study. Crucially, the template is \emph{prefilled
into the assistant turn} of the victim, not appended to the user
message --- the mask-fill objective then forces the diffusion victim to
complete the harmful specifics in parallel.

\lstinputlisting[style=appendixprompt]{tables/prompts/attacker.txt}

\subsection{Prompt for Scorer LLM}\label{appd:prompt-scorer}
The scorer prompt is a verbatim transcription of HarmBench's
\citep{harmbench} ``rule-based'' classifier system prompt, used here to
turn a victim generation into a reward $\rho \in \{0,1\}$ for the UCB
update and the success / fallback branch. We deliberately reuse the
upstream prompt unchanged so that our scorer signal is calibrated
against published benchmarks; the same prompt is also used by the
Bedrock judge in the headline tables (Sec.~\ref{appd:eval-metrics}).

\lstinputlisting[style=appendixprompt]{tables/prompts/scorer.txt}

\subsection{Prompt for Drafter LLM}\label{appd:prompt-drafter}
The drafter section actually contains two prompts that together form the
fallback path used when the main attacker template scores below
$\rho_{\text{fb}} = 0.7$. (A) A short \emph{rephraser} prompt converts the
goal into a more concrete, on-topic sentence; the rephrased goal is then
sent to the weaker base model (Qwen3-4B-Base) which drafts an uncensored
malicious response. (B) A \emph{Bedrock UNSAFE-span tagger} prompt
(Qwen3-235B) takes that draft and inserts \texttt{[UNSAFE]\dots[/UNSAFE]}
markers around the concrete operational spans; \method then replaces each
tagged span with a \verb|<mask:N>| of equal token length, producing the
next iteration's mask template. The combined fallback path lets \method
recover from cases where the attacker LLM's first scaffold is too coarse
to elicit harmful content from the victim.

\lstinputlisting[style=appendixprompt]{tables/prompts/drafter.txt}

\subsection{Prompt for Summarizer LLM}\label{appd:prompt-summarizer}
Whenever a Stage~B iteration improves on the previous best reward for the
same goal, the summarizer is invoked with both the previous and the
current $(template, schema, reward)$ tuple and asked to distil the
improvement into one new reusable pattern. The output is a tightly
formatted JSON object with three keys: \texttt{improvement\_summary} (free
text), \texttt{new\_pattern} (the structural schema with
\texttt{structure\_type}, \texttt{slot\_count}, \texttt{slot\_roles}, etc.
that becomes the canonical hash key in the registry), and
\texttt{new\_template} (a reusable goal-agnostic scaffold using
\texttt{<mask\_token>} as the placeholder). The summarizer is the
mechanism by which Stage~B's online attacks accumulate into a persistent
library --- on every improvement, a new pattern is hashed and inserted
into the registry, becoming available to all subsequent goals through
the UCB pool.

\lstinputlisting[style=appendixprompt]{tables/prompts/summarizer.txt}